% Generated draft: requires mathematical and source review.
\documentclass[11pt]{article}
\usepackage[margin=25mm]{geometry}
\usepackage{fontspec}
\setmainfont{Latin Modern Roman}
\usepackage{amsmath,amssymb,mathtools,unicode-math}
\setmathfont{Latin Modern Math}
\usepackage{enumitem,needspace,xurl,hyperref}
\setlength{\parindent}{0pt}
\setlength{\parskip}{5pt}
\setlength{\emergencystretch}{3em}
\newcommand{\N}{\mathbb{N}}
\newcommand{\Z}{\mathbb{Z}}
\newcommand{\Q}{\mathbb{Q}}
\newcommand{\R}{\mathbb{R}}
\newcommand{\C}{\mathbb{C}}
\newcommand{\F}{\mathbb{F}}
\newcommand{\A}{\mathbb{A}}
\newcommand{\PP}{\mathbb{P}}
\newcommand{\E}{\mathbb{E}}
\newcommand{\eps}{\varepsilon}
\newcommand{\dd}{\,\mathrm{d}}
\begin{document}
\section*{UnsolvedMath ID: 10900014}
\textbf{Source problem code:} catalogue entry
\subsection{Definitions and mathematical statement}
A finite-volume hyperbolic $3$-manifold is a complete hyperbolic $3$-manifold of finite Riemannian volume. A manifold is \emph{virtually} endowed with a specified property if some connected finite-sheeted cover has that property.

In the usual strict semi-bundle interpretation, a semi-fiber is a properly embedded separating surface $S$ in a compact manifold $N$ such that cutting along $S$ gives two interval bundles,
\[
 N=N_1\cup_S N_2,
\]
with $S$ the horizontal boundary of each interval bundle and with any vertical boundary contained in $\partial N$. In the closed orientable case, these are the familiar twisted interval bundles over nonorientable surfaces. For a cusped hyperbolic manifold one uses a compact core and the corresponding boundary convention. This is the interval-bundle interpretation suggested by the supplied summary, but the excerpt does not itself define semi-fibered.

\textbf{Question (Agol).} Does every finite-volume hyperbolic $3$-manifold $M$ admit a connected finite-sheeted covering $\widehat M\to M$ for which a compact core $\widehat N$ contains a semi-fiber, with the interval-bundle decomposition described above? In the closed case $\widehat N=\widehat M$. The assertion concerns finite-volume manifolds, not only the closed case.

The intended convention must be confirmed: some terminology allows fiber surfaces as semi-fibers, whereas the supplied summary expressly distinguishes this problem from virtual fibering over the circle. With that distinction, merely finding a finite cover that is a surface bundle over $S^1$ does not supply the requested two-interval-bundle structure. The excerpt also leaves orientability assumptions and the precise boundary conventions for cusped manifolds unstated.
\subsection{Short English statement}
Does every finite-volume hyperbolic $3$-manifold have a finite cover admitting the intended semi-fiber structure, namely a decomposition into two interval bundles along their horizontal boundaries? The source's exact convention for semi-fibering must be checked.
\subsection{Sources}
\begin{itemize}
\item Local UnsolvedMath registry, permanent ID 10900014, source code catalogue entry; attributed to Agol, no proposal year supplied.
\item Delp, Hoffoss and Manning, \emph{Problems In Groups, Geometry, and Three-Manifolds} (2015), subsection 3.6 (Agol), PDF page 2; \url{https://arxiv.org/pdf/1512.04620#page=2} and \url{https://arxiv.org/abs/1512.04620}. The registry records source-TeX extraction checked against rendered PDF text and access on 2026-07-29.
\item The supplied research summary distinguishes virtual semi-fibering from the reported virtual-fibering theorem for finite-volume hyperbolic $3$-manifolds. It gives no separate primary citation or precise semi-fiber definition.
\end{itemize}
\end{document}
